% lab_template.tex — neutrale Vorlage für Laborberichte mit labstyle.sty
\documentclass[a4paper,11pt]{scrartcl}
\KOMAoptions{
  parskip=half,
  headings=small,
  headinclude=true,
  footinclude=true
}
\areaset{175mm}{270mm}

\usepackage{labstyle}

% ------------------------------------------
% Sprache
% ------------------------------------------
% pdfLaTeX:
\usepackage[ngerman]{babel}
% LuaLaTeX / XeLaTeX:
% \usepackage{polyglossia}
% \setmainlanguage{german}

% ------------------------------------------
% Zitate
% ------------------------------------------
\usepackage{csquotes}

% ------------------------------------------
% Beispiele für eigene Einheiten (optional)
% ------------------------------------------
\DeclareSIUnit{\rpm}{rpm}
\DeclareSIUnit{\newtonmeter}{\newton\metre}

% ------------------------------------------
% Metadaten (Titelblatt)
% ------------------------------------------
\renewcommand{\labtitle}{Praktikum X}
\renewcommand{\labsubtitle}{Kurzer Untertitel des Versuchs}
\renewcommand{\labauthor}{Name / Gruppe}
\renewcommand{\labdate}{\today}

\begin{document}
\makelabtitle

% ==========================================================
% Abschnitt: Kurzbeschreibung / Lernziele
% ==========================================================
\section*{Kurzbeschreibung}
% Kurze Einleitung zum Versuch:
% - Was wird untersucht?
% - Was ist das Ziel?
% - Welche Grössen sind Ein- und Ausgänge?

Hier kurz den Inhalt und Zweck des Praktikums beschreiben.

\medskip
\textbf{Lernziele}
\begin{itemize}
    \item Lernziel 1 (z.\,B. Modellbildung eines dynamischen Systems)
    \item Lernziel 2 (z.\,B. Entwurf eines Reglers)
    \item Lernziel 3 (z.\,B. Umsetzung in Simulink / auf Hardware)
    \item Lernziel 4 (z.\,B. Auswertung und Interpretation der Ergebnisse)
\end{itemize}

\begin{answerboxfillpage}
% Platz für Antworten / Notizen zu den Lernzielen (falls benötigt)
\end{answerboxfillpage}

% ==========================================================
% Abschnitt: Systembeschreibung / Differentialgleichung
% ==========================================================
\section{Systembeschreibung und Modellierung}

% Beispielabbildung (Foto / Skizze des Aufbaus)
\labfig[width=0.9\linewidth]{example_setup.png}{Versuchsaufbau (Beispielbild).}

% Kurze verbale Beschreibung des Systems:
% - Welche physikalischen Grössen sind beteiligt?
% - Welche Vereinfachungen werden getroffen?

Hier System in Worten beschreiben.

% Beispiel: allgemeine Differentialgleichung
\[ 
    % Beispielstruktur:
    % a_2 \ddot{x}(t) + a_1 \dot{x}(t) + a_0 x(t) = b_0 u(t)
\]

\begin{gentable}{ll}{Systemparameter (Beispieltabelle)}
    \textbf{Grösse} & \textbf{Wert} \\
    \midrule
    Parameter 1 & Wert / Einheit \\
    Parameter 2 & Wert / Einheit \\
    Parameter 3 & Wert / Einheit \\
\end{gentable}

\begin{aufgaben}
    \item Skizzieren Sie die angreifenden Kräfte / Ströme / Momente am System.
    \item Benennen Sie Systemeingang und Systemausgang.
    \item Stellen Sie die Bewegungsgleichung / Systemgleichung auf.
\end{aufgaben}

\begin{answerboxfillpage}
% Platz für Herleitungen / Skizzen
\end{answerboxfillpage}

% ==========================================================
% Abschnitt: Linearisierung und Übertragungsfunktion
% ==========================================================
\section{Linearisierung und Übertragungsfunktion}

\subsection*{Systemanalyse}

\begin{aufgaben}
    \item Linearisieren Sie die nichtlineare Differentialgleichung um einen gegebenen Arbeitspunkt.
    \item Bestimmen Sie die Übertragungsfunktion vom Eingang $U(s)$ zum Ausgang $Y(s)$.
    \item Vereinfachen Sie das Modell ggf. (z.\,B. Reibung vernachlässigen, Totzeiten, etc.).
    \item Bestimmen Sie Pol- und Nullstellen der Übertragungsfunktion.
    \item Charakterisieren Sie das System hinsichtlich Stabilität und statischem Verhalten.
\end{aufgaben}

\begin{answerboxfillpage}
% Rechnungen / Ergebnisse zur Übertragungsfunktion
\end{answerboxfillpage}

\begin{answerbox}[8]
% Kompakte Zusammenfassung (z.B. finale Übertragungsfunktion, Pole, Nullstellen)
\end{answerbox}

% ==========================================================
% Abschnitt: Reglerentwurf (z.B. PI-Regler, Root Locus)
% ==========================================================
\section{Reglerentwurf}

% Beispiel PI-Regler
\[
    C(s) = k_p \left( \frac{T_i s + 1}{T_i s} \right)
\]

\begin{aufgaben}
    \item Wählen Sie eine geeignete Nachstellzeit $T_i$ (z.\,B. über Vorgabe einer Nullstelle).
    \item Bestimmen Sie die Übertragungsfunktion des offenen Regelkreises $L(s)$.
    \item Normieren Sie $L(s)$ auf $k \tilde{L}(s)$.
    \item Skizzieren Sie die Wurzelortskurve:
    \begin{enumerate}[label*=\alph*), leftmargin=2em]
        \item Pol- und Nullstellen von $\tilde{L}(s)$ eintragen.
        \item Polüberschuss und Asymptoten-Schnittpunkt bestimmen.
        \item Reelle Achsenanteile der Wurzelortskurve identifizieren.
        \item Verzweigungspunkte (breakaway/break-in) bestimmen.
        \item Wurzelortskurve vervollständigen.
        \item Verstärkung $k$ bzw. $k_p$ bestimmen, z.\,B. für doppelte Polstellen.
    \end{enumerate}
\end{aufgaben}

\begin{answerboxfillpage}
% Root-Locus-Konstruktion, Berechnungen, gewählte Reglerparameter
\end{answerboxfillpage}

% ==========================================================
% Abschnitt: Implementation und Test
% ==========================================================
\section{Implementation und Test}

% Hinweis auf Simulink / Implementierung auf Hardware:
% - Logik für Enable/Reset
% - Begrenzung von Stellgrössen
% - Initialbedingungen

\begin{aufgaben}
    \item Implementieren Sie den entworfenen Regler im gegebenen \emph{Simulink}-Modell.
    \item Führen Sie Versuche mit unterschiedlichen Sollwerten / Störungen durch.
    \item Dokumentieren Sie typische Verläufe (z.B. Führungsverhalten, Störverhalten).
\end{aufgaben}

\begin{answerboxfillpage}
% Versuchsbeschreibung, Screenshots von Verläufen, Beobachtungen
\end{answerboxfillpage}

% ==========================================================
% Abschnitt: Auswertung und Diskussion
% ==========================================================
\section{Auswertung und Diskussion}

\begin{aufgaben}
    \item Vergleichen Sie Messung und Modell (z.B. Sprungantwort, Frequenzgang).
    \item Diskutieren Sie Abweichungen und mögliche Ursachen.
    \item Beurteilen Sie die Qualität der Regelung (Überschwingen, Einregelzeit, Stellgrössenbegrenzung).
\end{aufgaben}

\begin{answerboxruled}[8]
% Zusammenhängender Text für die Auswertung / Diskussion
\end{answerboxruled}

% ==========================================================
% Abschnitt: Anhang / Zusatzaufgaben
% ==========================================================
\section*{Zusatzaufgaben}

\begin{aufgaben}[Zusatzaufgaben]
    \item Diskutieren Sie Vor- und Nachteile der gewählten Modellvereinfachungen.
    \item Erweitern Sie das Modell um ein nichtlineares Element (z.B. Reibung, Sättigung).
\end{aufgaben}

\begin{answerboxruled}[5]
% Platz für weiterführende Gedanken / Erweiterungen
\end{answerboxruled}

% ==========================================================
% Literatur (Beispieleintrag, anpassen/erweitern)
% ==========================================================
\begin{thebibliography}{1}

    \bibitem{BeispielQuelle}
    A.~Autor,
    \newblock ``Titel des Artikels oder Buches,''
    \newblock {\em Zeitschrift / Verlag}, Jahr.

\end{thebibliography}

\end{document}
